% !TEX program = xelatex
% HaUI / SICT research presentation template (English)
% Compiler: XeLaTeX (Overleaf: Menu > Compiler > XeLaTeX)
\documentclass[aspectratio=169, 11pt, t]{beamer}
\usetheme{HaUI}

\usepackage{booktabs}
\usepackage{pgfplots}
\pgfplotsset{compat=1.17}

% ----------------------------------------------------------- Metadata -------
\title[Short Research Title]{Title of the Research Paper Goes Here}
\subtitle{Optional subtitle or tagline}
\author[N. V. A, T. T. B]{Nguyen Van A, Tran Thi B}
\institute{SICT, Hanoi University of Industry}
\date{Hanoi, October 2026}
\confname{Student Research Conference 2026}
\contactinfo{quang.do@fit-haui.edu.vn \quad|\quad sict.haui.edu.vn}

\begin{document}

% ------------------------------------------------ 1. Title -----------------
\begin{frame}[plain]
  \titlepage
\end{frame}

% ------------------------------------------------ 2. Outline ---------------
\begin{frame}{Outline}
  \hauisidephoto
  \tableofcontents[hideallsubsections]
\end{frame}

% ------------------------------------------------ 3. Section + 4. Content ---
\section{Introduction}

\begin{frame}{Background and Motivation}
  \begin{columns}[T]
    \begin{column}{0.55\textwidth}
      \begin{itemize}
        \item Background: briefly describe the field and the practical problem.
        \item Research gap:
          \begin{itemize}
            \item Limitations of existing methods.
            \item Lack of data or evaluation for the target setting.
          \end{itemize}
        \item \alert{Research question}: state it clearly and make it measurable.
      \end{itemize}
    \end{column}
    \begin{column}{0.4\textwidth}
      \begin{block}{Objective}
        Propose a new method that improves accuracy while reducing computational cost.
      \end{block}
      \begin{exampleblock}{Key Contributions}
        A model, a dataset, and experimental results.
      \end{exampleblock}
      \begin{alertblock}{Note}
        Replace this sample text with your own content.
      \end{alertblock}
    \end{column}
  \end{columns}
\end{frame}

\section{Proposed Method}

\begin{frame}{Workflow}
  \begin{enumerate}
    \item Collect and preprocess the data.
    \item Build the proposed model.
    \item Train and tune the hyperparameters.
    \item Evaluate with standard metrics: $F_1 = \dfrac{2PR}{P+R}$.
  \end{enumerate}
\end{frame}

% ------------------------------------------------ 6. Experiments & Results ---
\section{Experiments \& Results}

% Layout: bảng kết quả + nhận xét
\begin{frame}{Experimental Setup and Main Results}
  \begin{columns}[T]
    \begin{column}{0.5\textwidth}
      \textbf{Setup}
      \begin{itemize}
        \item Dataset: describe source, size, train/val/test split.
        \item Baselines: ResNet-50~\cite{he2016resnet}, ViT-B/16~\cite{dosovitskiy2021vit}.
        \item Optimizer: Adam~\cite{kingma2015adam}, lr $=10^{-4}$.
        \item Metrics: Accuracy, $F_1$, inference time.
      \end{itemize}
    \end{column}
    \begin{column}{0.47\textwidth}
      \begin{table}
        \caption{Comparison on the test set (\%).}
        \centering\small
        \begin{tabular}{lccc}
          \toprule
          Method & Acc. & $F_1$ & Time (ms) \\
          \midrule
          ResNet-50~\cite{he2016resnet}      & 86.4 & 85.1 & 12.3 \\
          ViT-B/16~\cite{dosovitskiy2021vit} & 88.9 & 87.6 & 21.8 \\
          \midrule
          \textbf{Ours}                      & \textbf{91.2} & \textbf{90.3} & 14.1 \\
          \bottomrule
        \end{tabular}
      \end{table}
      \begin{alertblock}{Takeaway}
        +2.3 points of accuracy over the best baseline at lower cost.
      \end{alertblock}
    \end{column}
  \end{columns}
\end{frame}

% Layout: biểu đồ + diễn giải
\begin{frame}{Performance Analysis}
  \begin{columns}[T]
    \begin{column}{0.58\textwidth}
      \begin{figure}
        \centering
        \begin{tikzpicture}
          \begin{axis}[
            width=\linewidth, height=5.2cm, ybar, bar width=9pt,
            ymin=80, ymax=94, ylabel={Score (\%)},
            symbolic x coords={Acc.,Prec.,Recall,$F_1$}, xtick=data,
            legend style={at={(0.5,1.05)},anchor=south,legend columns=-1,draw=none},
            ymajorgrids, grid style={haui-bg}, tick label style={font=\small},
            axis line style={haui-text!50}]
            \addplot[fill=haui-text!40,draw=none] coordinates {(Acc.,86.4)(Prec.,85.9)(Recall,84.4)($F_1$,85.1)};
            \addplot[fill=haui-blue,draw=none]    coordinates {(Acc.,88.9)(Prec.,88.2)(Recall,87.0)($F_1$,87.6)};
            \addplot[fill=haui-gold,draw=none]    coordinates {(Acc.,91.2)(Prec.,90.8)(Recall,89.8)($F_1$,90.3)};
            \legend{ResNet-50, ViT-B/16, Ours}
          \end{axis}
        \end{tikzpicture}
        \caption{Metric comparison across methods.}
      \end{figure}
    \end{column}
    \begin{column}{0.38\textwidth}
      \begin{itemize}
        \item Consistent gains on all four metrics.
        \item Largest improvement on \alert{Recall}: +2.8 over ViT.
        \item Ablation: removing module A drops $F_1$ by 1.9 points.
      \end{itemize}
      \begin{exampleblock}{Insight}
        Explain \emph{why} the method works, not only that it works.
      \end{exampleblock}
    \end{column}
  \end{columns}
\end{frame}

% ------------------------------------------------ 7. Conclusion & Future Work -
\section{Conclusion \& Future Work}

\begin{frame}{Conclusion and Future Work}
  \begin{columns}[T]
    \begin{column}{0.48\textwidth}
      \begin{block}{Conclusion}
        \begin{itemize}
          \item Proposed a method that reaches \alert{91.2\%} accuracy.
          \item Cuts inference time by 35\% versus ViT-B/16~\cite{dosovitskiy2021vit}.
          \item Released code and dataset for reproducibility.
        \end{itemize}
      \end{block}
    \end{column}
    \begin{column}{0.48\textwidth}
      \begin{exampleblock}{Future Work}
        \begin{enumerate}
          \item Evaluate on larger, multi-domain datasets.
          \item Model compression for edge deployment.
          \item Self-supervised pre-training with unlabeled data.
        \end{enumerate}
      \end{exampleblock}
    \end{column}
  \end{columns}
  \vspace{0.3cm}
  \begin{alertblock}{Limitations}
    Results are reported on a single dataset; generalization remains to be verified.
  \end{alertblock}
\end{frame}

% ------------------------------------------------ 8. References ---------------
% Dùng [allowframebreaks] nếu danh sách dài hơn một slide.
\begin{frame}[allowframebreaks]{References}
  \footnotesize
  \begin{thebibliography}{9}
    \bibitem{he2016resnet}
      K.~He, X.~Zhang, S.~Ren, J.~Sun.
      \newblock Deep residual learning for image recognition.
      \newblock In \emph{Proc. IEEE CVPR}, pp.~770--778, 2016.
    \bibitem{dosovitskiy2021vit}
      A.~Dosovitskiy \emph{et al.}
      \newblock An image is worth 16x16 words: Transformers for image recognition at scale.
      \newblock In \emph{Proc. ICLR}, 2021.
    \bibitem{kingma2015adam}
      D.~P.~Kingma, J.~Ba.
      \newblock Adam: A method for stochastic optimization.
      \newblock In \emph{Proc. ICLR}, 2015.
    \bibitem{vaswani2017attention}
      A.~Vaswani \emph{et al.}
      \newblock Attention is all you need.
      \newblock In \emph{Advances in NeurIPS}, vol.~30, 2017.
  \end{thebibliography}
\end{frame}

% ------------------------------------------------ 5. Thank you --------------
\hauithanksframe

\end{document}
